% AUTORIDAD DE FUENTE:
% expose_VIII_body.tex, líneas 2449--2495; índices de escaneo 321--323;
% folios de fuente 310--312. No hay diferencias materiales entre capas;
% los acentos bibliográficos siguen la autoridad impresa.

\medskip
\noindent\underline{Proposición 7.4. Sea $S$ un esquema reducido.}

\smallskip
\noindent a) \underline{Sea $P$ un esquema en grupos liso sobre $S$
cuyas fibras en los puntos maximales de $S$ son conexas. Entonces el
funtor natural}
\begin{equation}\tag{7.4.1}
 \mathrm{EXT}(P,\underline G_{mS})
 \longrightarrow
 \mathrm{TORSRIG}(P,e_P;\underline G_{mS})
\end{equation}
\underline{de la categoría de extensiones de $P$ por
$\underline G_{mS}$ a la categoría de torsors bajo
$\underline G_{mP}$ rigidificados a lo largo de la sección identidad
de $P$ es plenamente fiel. Si $S$ es normal y $P$ tiene fibras conexas,
un objeto $E$ del destino de (7.4.1) pertenece a la imagen esencial de
este funtor si y solo si, para todo punto maximal $\eta$ de $S$, el
torsor rigidificado $E_\eta$ sobre $P_\eta$ procede de una extensión de
$P_\eta$ por $\underline G_{m\eta}$.}

\smallskip
\noindent b) \underline{Sean $P,Q$ dos esquemas en grupos lisos sobre
$S$ cuyas fibras en los puntos maximales de $S$ son conexas. Entonces
el funtor natural}
\begin{equation}\tag{7.4.2}
 \mathrm{BIEXT}(P,Q;\underline G_{mS})
 \longrightarrow
 \mathrm{TORSBIRIG}(P,Q;\underline G_{mS})
\end{equation}
\underline{de la categoría de biextensiones de $(P,Q)$ por
$\underline G_{mS}$ a la categoría de torsors bajo
$\underline G_{m,P\times_S Q}$ birrigidificados a lo largo de la
pareja de secciones identidad de $P$ y $Q$ es plenamente fiel. Si $S$
es normal y $P$ o $Q$ tiene fibras conexas, un objeto $E$ del destino
de (7.4.2) pertenece a la imagen esencial del funtor si y solo si, para
todo punto maximal $\eta$ de $S$, $E_\eta$ procede de una biextensión de
$(P_\eta,Q_\eta)$ por $\underline G_{m\eta}$.}

Demostremos (a); la demostración esencialmente idéntica de (b) se deja
al lector. Para la plena fidelidad, obsérvese que 4.1 sigue siendo
válido si $X$ e $Y$ son esquemas planos sobre un esquema reducido $S$
cuyas fibras en los puntos maximales satisfacen las condiciones de 4.1.
En efecto, la familia de fibras maximales de $X\times_S Y$ —que es
plano sobre el esquema reducido $S$— es esquemáticamente densa en
$X\times_S Y$. Por tanto, un morfismo de $X\times_S Y$ al esquema
separado $\underline G_{mS}$ sobre $S$ es trivial en cuanto lo sea
sobre las fibras maximales. La plena fidelidad de (7.4.1) se demuestra
entonces como en 4.2, y la caracterización de su imagen esencial dada
allí sigue siendo válida.

Queda por demostrar que, si el torsor rigidificado $E$ sobre $P$ es tal
que los $E_\eta$ proceden de extensiones, entonces también $E$ procede
de una extensión, suponiendo que $S$ es normal y $P$ tiene fibras
conexas. Esta última condición implica que $P$ es de presentación
finita sobre $S$ (SGA 3 $\mathrm{VI}_B$ 5.5). La plena fidelidad muestra
también que la cuestión es local sobre $S$, de modo que podemos suponer
$S$ afín y local. El argumento usual de paso al límite, utilizando el
teorema de Nagata sobre la finitud de la clausura integral de una
$\mathbb Z$-álgebra reducida y finitamente generada (EGA IV 7.8), nos
reduce además al caso en que $S$ es noetheriano.

%%%%%%%%%%  scan idx 322  |  folio 311  %%%%%%%%%%

Entonces 7.1(a) muestra que existe un abierto $U\subset S$ cuyo
complemento tiene codimensión al menos $2$ y tal que $E_U$ procede de
una extensión de $P_U$ por $\underline G_{mU}$. En otras palabras, el
isomorfismo buscado
\[
 \pi^*(E)\simeq
 \operatorname{pr}_1^*(E)\operatorname{pr}_2^*(E)
\]
puede construirse sobre $(P\times_S P)_U$. Este es un abierto del
esquema normal $P\times_S P$ cuyo complemento tiene codimensión al menos
$2$. Por tanto, el isomorfismo sobre $(P\times_S P)_U$ se prolonga a
todo $P\times_S P$, lo que concluye la demostración.

%%%%%%%%%%  scan idx 323  |  folio 312  %%%%%%%%%%

\medskip
\noindent\underline{Corolario 7.5. Sea $S$ un esquema normal y sean $P$
y $Q$ dos esquemas en grupos lisos sobre $S$, con fibras conexas en los
puntos maximales, y supongamos que $P$ o $Q$ tiene fibras conexas.
Supongamos además que, para todo punto maximal $\eta$ de $S$, $P_\eta$
o $Q_\eta$ es una extensión de un esquema abeliano por un toro (o, más
generalmente, admite una serie de composición cuyos factores son
esquemas abelianos, toros o grupos $\underline G_a$). Entonces el funtor
(7.4.2) es una equivalencia de categorías.}

Es una consecuencia inmediata de 7.4(b) y 4.7.

\bigskip
\noindent\underline{BIBLIOGRAFÍA}

\begin{enumerate}
\item[{[1]}] M. Demazure, J. Giraud, M. Raynaud,
  \emph{Schémas abéliens}, Séminaire Orsay 1967/68, de próxima
  aparición.
\item[{[2]}] M. Raynaud, \emph{Faisceaux amples sur les schémas en
  groupes et les espaces homogènes}, tesis, París, 1968 (de próxima
  aparición en Lecture Notes, Springer).
\item[{[3]}] M. Rosenlicht, \emph{Some basic theorems on algebraic
  groups}, Amer. Jour. Math. \textbf{78}, n.º~2 (1956), pp.~401--443.
\item[{[3bis]}] M. Rosenlicht, \emph{Some rationality questions on
  algebraic groups}, Annali di Matematica, serie 4, n.º~43 (1957),
  pp.~25--50.
\item[{[4]}] J.-P. Serre, \emph{Corps locaux}, Act. Sci. Ind. 1296
  (1962), Hermann, París.
\end{enumerate}
